\xiaosi

\subsubsection{仿真环境}

本文使用Highway-env\cite{leurent2018highway}作为实验平台，选取4个难度递进的驾驶场景：高速公路巡航（highway-v0，40秒）、匝道合流（merge-v0，约40秒）、无保护交叉路口（intersection-v0，13秒）、环岛通行（roundabout-v0，11秒）。所有场景统一配置为5个离散元动作（左变道、保持、右变道、加速、减速），仿真频率5Hz。

\subsubsection{观测空间}

所有环境使用相同的运动学观测：每辆车7个特征（存在标志、相对位置$x$/$y$、相对速度$v_x$/$v_y$、朝向$\cos h$/$\sin h$），观测5辆车（自车+最近4辆），归一化后展平为35维向量。碰撞奖励设为$-5$以强化安全约束。

\subsubsection{对比模型}

实验对比6种模型：（1）随机策略，作为性能下界；（2）MLP（13,189参数），两层全连接网络；（3）LSTM（26,181参数），单层64维；（4）GRU（19,717参数），单层64维；（5）全连接LTC（7,152参数），与NCP参数量相同但采用全连接布线；（6）NCP（7,152参数），本文方法。所有模型使用相同的Double DQN训练配置。

\subsubsection{训练配置}

学习率$3 \times 10^{-4}$（Adam优化器），折扣因子$\gamma = 0.99$，批量大小64，经验回放容量50,000，目标网络每500步硬更新，$\varepsilon$从1.0线性衰减至0.05（10,000步内），总训练50,000步，梯度裁剪阈值10.0。循环模型使用长度8的序列采样。每组实验使用3个随机种子（0、42、123）确保统计显著性。全部实验在两张NVIDIA A10 GPU上并行执行，72组基线实验约2小时完成。
